\section{Protocollo operativo per similarit\`a e analisi del ciclo}
\label{sec:operational-experiment-protocol}
Questa sezione precisa come vengono prodotte e lette le osservazioni,
senza introdurre fattori aggiuntivi nel piano a cinque configurazioni.

\subsection{Dalle descrizioni alle etichette cosine}
Per ciascun dataset si conserva la baseline top-down e si raggruppano
gli LLG secondo il riferimento al padre. La ricostruzione riceve soltanto
le descrizioni dei figli non vuoti: non vede il padre, la ground truth
o il documento per riprodurli. Qwen restituisce una descrizione HLG$'$
in JSON; richieste e retry di ricostruzione non riutilizzano risposte
precedentemente scartate. I rami non ricostruibili sono registrati,
senza trasformare un dato mancante in similarit\`a zero.

MPNet codifica separatamente la descrizione ricostruita e quelle degli
HLG originali correnti, con vettori normalizzati. Il codice usa CPU,
cache delle rappresentazioni e nessun preprocessing lessicale aggiuntivo;
controlla testi vuoti, vettori non validi e superamento del limite di
token prima di ammettere il confronto. Per vettori unitari il prodotto
scalare coincide con la cosine:
\begin{equation}
u_i=\frac{e(HLG'_i)}{\|e(HLG'_i)\|},\quad
v_j=\frac{e(HLG_j)}{\|e(HLG_j)\|},\qquad s_{ij}=u_i^\top v_j.
\end{equation}
Ogni HLG$'_i$ viene confrontato con tutti gli HLG originali dello
stesso progetto: $j=i$ indica il proprio padre; $j\ne i$ un altro
padre. Non si usano confronti HLG$'$--HLG$'$ e non si incrociano dataset.
Con $\tau=0{,}68$, il confronto col padre produce
\texttt{EQUIVALENT} quando $s_{ii}\ge\tau$ e
\texttt{NOT\_EQUIVALENT} altrimenti. I confronti fuori diagonale
sopra soglia possono diventare candidati duplicati soltanto per HLG
dello stesso attore e con le disponibilit\`a richieste dal codice.
L'uguaglianza dell'attore usa il nome normalizzato per spazi e maiuscole;
non \`e una verifica semantica automatica dei ruoli.

Gli evaluator ricevono etichette fisse e coppie candidate testuali,
non i punteggi da riclassificare. La somiglianza numerica serve a
individuare confronti da verificare, mentre la documentazione determina
supporto e responsabilit\`a della funzionalit\`a. La soglia \`e comune
ai dataset e non viene adattata in base alla ground truth. Nello studio
esplorativo si ricalcolano i conteggi sopra soglia sulle stesse coppie
per diversi valori ipotizzati; non si eseguono altrettante pipeline
iterative. Il raggruppamento diretto di quel run, la sua copertura e
i valori osservati restano quelli descritti nella
Sezione~\ref{sec:cosine-tuning-all-datasets}.

\subsection{Esecuzione dei confronti di prompting e temperatura}
Ogni run riparte dalla baseline congelata, anzich\'e dall'output della
configurazione precedente. Il confronto P0--P1--P4 varia soltanto
il prefisso di scenari dimostrativi a temperatura 0,2; il confronto
T0--P4--T4 mantiene quattro scenari e varia la temperatura fra
0, 0,2 e 0,4. P4 \`e una sola condizione condivisa. Ricostruzione,
generazione e critic della riparazione mantengono le proprie impostazioni.
La verifica iniziale del backend controlla l'accettazione delle
temperature richieste; non dimostra identit\`a delle risposte a temperatura zero.
Una modifica di prompt o codice richiede una suite distinta e non viene
applicata retroattivamente ai dati storici. Le analisi del processo
di seguito riguardano P0, P1 e P4 della v111, tutte terminate sui dodici
dataset; i risultati parziali delle temperature non vi entrano.

\subsection{Unit\`a di conteggio per i singoli LLM e per il ciclo}
Si distinguono tre livelli: la risposta globale di un modello a una
iterazione; il tentativo tecnico usato per ottenerla; la decisione
di quel modello su un ramo all'interno della risposta accettata.
Un tentativo scartato non fornisce voti parzialmente utilizzabili.
La percentuale di azioni di un LLM ha come denominatore le sue decisioni
accettate, non il numero di richieste o di dataset. Le frequenze sono
quindi pesate dall'esposizione: run lunghe e gerarchie grandi contribuiscono
pi\`u delle run brevi. Non costituiscono accuratezza dei singoli modelli.

L'accordo confronta azione e, per una rimozione duplicata, superstite.
L'accordo a coppie usa solo rami per cui entrambi i modelli hanno una
risposta valida. Il quorum \`e misurato separatamente dall'applicazione
della proposta. Un voto di riscrittura non equivale a una riparazione
riuscita; \texttt{HLG\_LLG\_HANDOFF\_COMPLETED} registra il completamento
del passaggio della collezione, non il conteggio dei goal effettivamente
modificati. Pi\`u rami possono riferirsi alla medesima riparazione batch.

La convergenza \`e l'arresto per conferma di tutti i rami, mentre il
limite di 25 iterazioni \`e un diverso motivo di arresto. Per verificare
le modifiche si confrontano le collezioni HLG--LLG in ingresso e in
uscita della stessa iterazione. Le ripetizioni esatte di stato sono
rilevate mediante impronta del contenuto serializzato, preservando
l'ordine delle collezioni. Gli ID di ramo sono locali allo snapshot:
non vengono usati per sostenere che lo stesso goal oscilli dopo una
rigenerazione che potrebbe aver cambiato cardinalit\`a o ordine.
